\clearpage\pagestyle{otMain}\setlength{\otBlocEspacement}{2.5mm}\fontsize{11}{15}\selectfont\corrigesChapitre{18}{Solides, patrons et espace}
\begin{blocLivre}{cor-18-1}\begin{corrige}{1}{Nommer les solides}\demarche{Observer les faces et les surfaces courbes.}
Dans l’ordre : $\res{\text{cube}}$, $\res{\text{pavé}}$, $\res{\text{cylindre}}$, $\res{\text{pyramide}}$, $\res{\text{cône}}$ et $\res{\text{boule}}$.
\label{fig-cor-18-1}\figureGeo[corrige]{sol-1}
\end{corrige}\end{blocLivre}
\begin{blocLivre}{cor-18-2}\begin{corrige}{2}{La carte d’identité du cube}\demarche{Compter sur la perspective, faces cachées comprises.}
Le cube a $\res6$ faces $\res{\text{carrées}}$, $\res{12}$ arêtes et $\res8$ sommets.
\label{fig-cor-18-2}\figureGeo[corrige]{sol-2}
\end{corrige}\end{blocLivre}
\begin{blocLivre}{cor-18-3}\begin{corrige}{3}{Décrire un pavé}\demarche{S’appuyer sur la comparaison avec le cube.}
Le pavé a $\res6$ faces $\res{\text{rectangulaires}}$, 12 arêtes et 8 sommets. Les faces opposées sont $\res{\text{identiques}}$ deux à deux.
\label{fig-cor-18-3}\figureGeo[corrige]{sol-3}
\end{corrige}\end{blocLivre}
\begin{blocLivre}{cor-18-4}\begin{corrige}{4}{Les faces de la pyramide}\demarche{Séparer la base des faces latérales.}
Elle a une base $\res{\text{carrée}}$ et $\res4$ faces latérales $\res{\text{triangulaires}}$, soit 5 faces en tout.
\label{fig-cor-18-4}\figureGeo[corrige]{sol-4}
\end{corrige}\end{blocLivre}
\begin{blocLivre}{cor-18-5}\begin{corrige}{5}{Le prisme droit}\demarche{Chercher d’abord les deux bases jumelles.}
Ses deux bases sont des $\res{\text{triangles}}$ superposables ; les $\res3$ autres faces sont des $\res{\text{rectangles}}$. C’est un prisme droit à base triangulaire.
\label{fig-cor-18-5}\figureGeo[corrige]{sol-5}
\end{corrige}\end{blocLivre}
\begin{blocLivre}{cor-18-6}\begin{corrige}{6}{Patron ou pas ?}\demarche{Plier mentalement face par face.}
La croix et la disposition en T se plient en cube : $\res{\text{patrons}}$. La disposition où deux carrés sont du même côté d’une ligne de quatre fait se superposer deux faces : $\res{\text{pas un patron}}$.
\label{fig-cor-18-6}\figureGeo[corrige]{sol-6}
\end{corrige}\end{blocLivre}
\begin{blocLivre}{cor-18-7}\begin{corrige}{7}{Les faces opposées}\demarche{Sauter une face sur le patron.}
En pliant, la face opposée à A est celle qui en est séparée par une face sur le patron : c’est la face marquée $\res C$. Deux faces voisines sur le patron ne peuvent pas être opposées.
\label{fig-cor-18-7}\figureGeo[corrige]{sol-7}
\end{corrige}\end{blocLivre}
\begin{blocLivre}{cor-18-8}\begin{corrige}{8}{Compléter le patron}\demarche{Tester chaque position par pliage mental.}
Placé à droite du deuxième carré de la colonne, le sixième carré ferme le cube : cette position $\res{\text{convient}}$ ; elle figure en vert. Collé sous le carré de gauche, il formerait un bloc de quatre carrés : deux faces se superposeraient au pliage, position $\res{\text{interdite}}$.
\label{fig-cor-18-8}\figureGeo[corrige]{sol-8}
\end{corrige}\end{blocLivre}
\begin{blocLivre}{cor-18-9}\begin{corrige}{9}{Les cubes cachés}\demarche{Compter étage par étage, du bas vers le haut.}
Étage du bas : $3\times2=6$ cubes. Étage du milieu : 3 cubes. Étage du haut : 1 cube. Total : $6+3+1=\res{10}$ cubes ; les cubes porteurs cachés sont comptés.
\label{fig-cor-18-9}\figureGeo[corrige]{sol-9}
\end{corrige}\end{blocLivre}
\begin{blocLivre}{cor-18-10}\begin{corrige}{10}{La bonne vue}\demarche{Imaginer l’assemblage depuis chaque position.}
La vue de face montre les hauteurs des colonnes : vignette $\res2$. La vue de dessus montre les cases occupées au sol : vignette $\res3$. La vignette 1 ne correspond à aucun point de vue.
\label{fig-cor-18-10}\figureGeo[corrige]{sol-10}
\end{corrige}\end{blocLivre}
\begin{blocLivre}{cor-18-11}\begin{corrige}{11}{Construire un patron}\demarche{Tous les carrés ont le même côté de 3 cm.}
On dessine une croix de six carrés de $\res{3\ \mathrm{cm}}$ de côté, à la règle et à l’équerre ; les carrés partagent leurs côtés exactement. Le patron, plié, referme le cube : le tracé figure en vert.
\label{fig-cor-18-11}\figureGeo[corrige]{sol-11}
\end{corrige}\end{blocLivre}
\begin{blocLivre}{cor-18-12}\begin{corrige}{12}{Agrandir l’assemblage}\demarche{Calculer le total visé avant de soustraire.}
Un cube de côté 2 contient $2\times2\times2=8$ cubes : il manque $8-7=\res1$ cube. Un cube de côté 3 en contient $3\times3\times3=27$ : il manquerait $27-7=\res{20}$ cubes.
\label{fig-cor-18-12}\figureGeo[corrige]{sol-12}
\end{corrige}\end{blocLivre}
\begin{blocLivre}{cor-18-13}\begin{corrige}{13}{Coder un chemin}\demarche{Suivre l’orientation du pion, pas celle du lecteur.}
Un codage correct est : avancer de 3 cases ; pivoter d’un quart de tour à droite ; avancer de 1 case ; pivoter d’un quart de tour à droite ; avancer de 3 cases. L’arrivée est bien la case but, une case $\res{\text{sous le départ}}$ après le retour.
\label{fig-cor-18-13}\figureGeo[corrige]{sol-13}
\end{corrige}\end{blocLivre}
\begin{blocLivre}{cor-18-14}\begin{corrige}{14}{Exécuter un programme}\demarche{Mettre à jour l’orientation après chaque quart de tour.}
Le pion va 2 cases vers la droite, monte de 2 cases, puis revient de 1 case vers la gauche : il arrive 1 case à droite et 2 cases au-dessus du départ, tourné vers la $\res{\text{gauche de la feuille}}$. Le chemin figure en vert.
\label{fig-cor-18-14}\figureGeo[corrige]{sol-14}
\end{corrige}\end{blocLivre}
\begin{blocLivre}{cor-18-15}\begin{corrige}{15}{Le dé à compléter}\demarche{Associer les paires qui totalisent 7, puis les placer.}
La face opposée à 3 porte $7-3=\res4$ ; celle opposée à 2 porte $\res5$ ; celle opposée à 1 porte $\res6$. Sur le patron, chaque paire opposée est séparée par une face ; les valeurs figurent en vert.
\label{fig-cor-18-15}\figureGeo[corrige]{sol-15}
\end{corrige}\end{blocLivre}
